\section{CNN 的完整流程：从特征提取到分类}

\subsection{完整架构概览}

一个完整的 CNN 分类网络可以分为两个阶段：

\begin{figure}[H]
\centering
\includegraphics[width=0.9\textwidth]{slides-images/slide-51.jpg}
\caption{CNN 的完整流程：前半部分是\textbf{特征学习}（卷积 + ReLU + 池化的重复），后半部分是\textbf{分类}（展平 + 全连接 + Softmax）。\protect\footnotemark}
\end{figure}
\footnotetext{Slides 第51页。}

\begin{enumerate}
    \item \textbf{特征学习阶段}：
    \begin{itemize}
        \item 通过\textbf{卷积}学习输入图像中的特征
        \item 通过\textbf{非线性激活}引入非线性（因为真实世界的数据是非线性的）
        \item 通过\textbf{池化}降低维度并保留空间不变性
        \item 重复上述过程若干次
    \end{itemize}
    \item \textbf{分类阶段}：
    \begin{itemize}
        \item 将最后的特征图\textbf{展平}（flatten）为一维向量
        \item 通过\textbf{全连接层}进行分类
        \item 使用 \textbf{Softmax} 输出各类别的概率
    \end{itemize}
\end{enumerate}

\subsection{分类头：Softmax 输出}

CNN 的卷积和池化层输出高层特征表示后，通过全连接层将其映射到各类别，最后使用 Softmax 函数输出概率分布：

\begin{figure}[H]
\centering
\includegraphics[width=0.9\textwidth]{slides-images/slide-52.jpg}
\caption{CNN 的分类头：CONV 和 POOL 层输出高层特征，展平后通过全连接层，最终 Softmax 输出每个类别的概率。\protect\footnotemark}
\end{figure}
\footnotetext{Slides 第52页。}

$$\text{softmax}(y_i) = \frac{e^{y_i}}{\sum_j e^{y_j}}$$

其中 $y_i$ 是全连接层输出的第 $i$ 个类别的原始分数（logit），Softmax 将其转换为概率值，所有类别的概率之和为 1。

\subsection{代码实现：TensorFlow 与 PyTorch}

\begin{figure}[H]
\centering
\includegraphics[width=0.9\textwidth]{slides-images/slide-53.jpg}
\caption{TensorFlow/Keras 实现 CNN：两个卷积层（32 和 64 个滤波器）+ 展平 + 全连接层 + 10 类 Softmax 输出。\protect\footnotemark}
\end{figure}
\footnotetext{Slides 第53页。}

\begin{lstlisting}[caption={TensorFlow/Keras CNN 实现}]
import tensorflow as tf

def generate_model():
    model = tf.keras.Sequential([
        # First convolutional layer
        tf.keras.layers.Conv2D(32, filter_size=3, activation='relu'),
        tf.keras.layers.MaxPool2D(pool_size=2, strides=2),

        # Second convolutional layer
        tf.keras.layers.Conv2D(64, filter_size=3, activation='relu'),
        tf.keras.layers.MaxPool2D(pool_size=2, strides=2),

        # Fully connected classifier
        tf.keras.layers.Flatten(),
        tf.keras.layers.Dense(1024, activation='relu'),
        tf.keras.layers.Dense(10, activation='softmax')  # 10 outputs
    ])
    return model
\end{lstlisting}

\begin{lstlisting}[caption={PyTorch CNN 实现}]
import torch
import torch.nn as nn

def generate_model():
    model = nn.Sequential(
        # First and second convolutional layer
        nn.Conv2d(in_channels=3, out_channels=32, kernel_size=3),
        nn.ReLU(),
        nn.MaxPool2d(kernel_size=2, stride=2),

        nn.Conv2d(in_channels=32, out_channels=64, kernel_size=3),
        nn.ReLU(),
        nn.MaxPool2d(kernel_size=2, stride=2),

        # Fully connected classifier
        nn.Flatten(),
        nn.Linear(64*6*6, 1024),  # flattened dim after 2 conv layers
        nn.ReLU(),
        nn.Linear(1024, 10),      # 10 outputs
    )
    return model
\end{lstlisting}

\begin{knowledgebox}{超参数选择的"艺术与科学"}
Amini 坦言，CNN 超参数的选择（滤波器数量、大小、层数等）"更多是一种艺术而非科学"（more art than science），但也有一些经验法则：
\begin{itemize}
    \item 滤波器大小通常较小（$3 \times 3$ 或 $5 \times 5$），因为局部特征不需要很大的感受野
    \item 滤波器数量通常\textbf{逐层递增}（如 $32 \rightarrow 64 \rightarrow 128$），因为更高层需要表示更多种类的复合特征
    \item 空间尺寸通过池化\textbf{逐层递减}，与滤波器数量的递增形成互补
\end{itemize}
\end{knowledgebox}

\subsection{本章小结}

完整的 CNN 分类网络分为特征学习和分类两个阶段。特征学习通过卷积+ReLU+池化的堆叠逐层提取抽象特征；分类阶段通过展平和全连接层将特征映射为类别概率。TensorFlow 和 PyTorch 都提供了简洁的 API 来构建 CNN。

\newpage
